\section{Compound Sparsity Formulation}
\label{app:formulation-details}

\paragraph{Attention head pruning.}
Heads are ranked by Frobenius norm importance $w_j^{(\ell)} = \|\mathbf{A}_j^{(\ell)}\|_F$ \cite{voita-acl2019, michel-neurips2019}; a binary mask $m_j^{(\ell)} \in \{0,1\}$ retains the top-$(1{-}s)\cdot H$ heads:
\begin{equation}
    \tilde{\mathbf{o}}_j^{(\ell)} = m_j^{(\ell)} \cdot \mathbf{o}_j^{(\ell)}.
    \label{eq:head-mask}
\end{equation}

\paragraph{KV cache eviction.}
A position mask $\mathbf{m}_k \in \{0,1\}^T$ retains the top-$(1{-}s)\cdot T$ positions (BOS always preserved):
\begin{equation}
    \tilde{a}_{ij} = a_{ij} \cdot m_{k,j}, \quad \forall\, i, j.
    \label{eq:kv-mask}
\end{equation}
Under importance-based eviction, each position $j$ is scored by its cumulative attention mass with position normalization:
\begin{equation}
    \mathrm{imp}(j) = \frac{\sum_{\ell, h, i} a_{ij}^{(\ell,h)}}{T - j},
    \label{eq:kv-importance}
\end{equation}
where the denominator compensates for causal bias---earlier positions are attended to by more queries \cite{thinkv-iclr2026, defensivekv-iclr2026}. Under random eviction, positions (excluding BOS) are evicted uniformly at random.

\paragraph{Compound sparsity.}
Applying both mechanisms simultaneously yields compound sparsity. The two axes interact: head pruning redistributes attention mass, altering the importance ranking for KV eviction. Prior work \cite{spark-aaai2026} reports super-additive joint effects but treats interaction magnitude as a fixed property of $s$.

\paragraph{Attention sparsity metric.}
For each query position $i$ and head $h$, we compute the fraction of key positions receiving negligible attention:
\begin{equation}
    \sigma_i^{(h)} = \frac{1}{T} \sum_{j=1}^{T} \mathbf{1}\!\left[a_{ij}^{(h)} < 0.01 \cdot \max_k a_{ik}^{(h)}\right].
    \label{eq:attn-sparsity}
\end{equation}
The per-position sparsity $\sigma_i = \frac{1}{H}\sum_h \sigma_i^{(h)}$ averages across all heads. Tokens are partitioned into high-, mid-, and low-sparsity phases by the tercile boundaries of $\{\sigma_i\}$; segments shorter than 10 tokens are merged into adjacent phases.

\section{Multiplicative Null Verification}
\label{app:multiplicative-null}

We verified that replacing the additive null $r_b = r_a + r_k - 1$ (Eq.~\ref{eq:ii}) with the multiplicative null $r_b = r_a \cdot r_k$ does not change qualitative conclusions. Under the multiplicative formulation, the interaction index becomes $\mathrm{II}_{\times} = r_b - r_a \cdot r_k$. Notably, $\mathrm{II}_{\times} < 0$ across all phases and sparsity levels: the compound interaction is \emph{sub-multiplicative} everywhere---actual compound degradation is less severe than the multiplicative prediction. $G_{\times}$ therefore quantifies the ratio of sub-multiplicative deviations between phases, not super-multiplicative amplification: a value of $G_{\times}{=}30$ means the low-sparsity phase deviates from the multiplicative null by $30{\times}$ more than the high-sparsity phase. Under importance-based eviction, $G_{\times}{=}21.1$--$36.2{\times}$ (DeepSeek) and $43.9$--$64.0{\times}$ (Qwen) across $s \in \{0.15, 0.20, 0.25, 0.30\}$, yielding a combined range of $21$--$64{\times}$. Although compound effects do not exceed the multiplicative prediction in any phase, the low-sparsity phase's deviation from the multiplicative null is $21$--$64{\times}$ larger than the high-sparsity phase's, indicating that the phase-gated severity gradient is preserved under the multiplicative null. These $n{=}500$ exact values are more conservative than the preliminary $n{=}50$ calibrated estimates ($\approx$51--183${\times}$), reflecting reduced variance at full sample size. Values are computed using aggregate per-phase PPL ratios; per-problem exact computation is left for future work. Random eviction remains flat ($G_{\times} \approx 1$). The crossover at $s{=}0.25$ is preserved. We adopt the additive formulation in the main text as the more conservative baseline: it assumes zero interaction (compound loss equals the sum of individual losses), so rejection ($\mathrm{II} > 0$) establishes super-additivity regardless of null specification. The qualitative conclusions---phase-modulated severity, strategy asymmetry, and crossover---are preserved under the multiplicative null.

% Bootstrap and permutation details moved to main text \S\ref{sec:stat-validation}
